\documentclass[12pt]{article}
% Adaptado de C:/GitHub/Mestrado/Template/Template.tex.
\usepackage[T1]{fontenc}
\usepackage[utf8]{inputenc}
\usepackage[brazilian]{babel}
\usepackage{graphicx,booktabs,array,amsmath,amssymb,microtype}
\usepackage{geometry}
\usepackage{pdflscape,caption}
\geometry{a4paper,margin=2.5cm}
\usepackage{parskip,csquotes}
\interfootnotelinepenalty=10000
\usepackage[backend=biber,style=abnt]{biblatex}
\addbibresource{references.bib}
\usepackage{hyperref}
\hypersetup{colorlinks=true,urlcolor=blue,linkcolor=blue,citecolor=blue}
\newcommand{\instituicao}{UNIVERSIDADE FEDERAL DE SANTA CATARINA}
\newcommand{\departamento}{Programa de Pós-Graduação em Economia}
\newcommand{\curso}{Mestrado}
\newcommand{\disciplina}{Desigualdade, Diversidade e Redes}
\newcommand{\titulo}{Rede de emissões de CO\textsubscript{2} no Brasil}
\newcommand{\autor}{Gabriel Pincelli Cintra}
\newcommand{\local}{Florianópolis}
\newcommand{\dataentrega}{18 de outubro de 2026}
\newcommand{\imprimirtitulodois}{
\begin{center}
  {\Large\bfseries\MakeUppercase{\instituicao}\par}
  \vspace{0.2em}{\large \departamento\par}
  \vspace{0.2em}{\normalsize\textbf{\disciplina}\ (\curso)\par}
  \vspace{1.4em}{\LARGE\bfseries \titulo\par}
  \vspace{1.5em}
\end{center}
\noindent\begin{minipage}{0.48\textwidth}
  \raggedright\normalsize \autor
\end{minipage}\hfill
\begin{minipage}{0.48\textwidth}
  \raggedleft\normalsize \local, \dataentrega
\end{minipage}
\vspace{1.0em}\hrule\vspace{1.0em}}
\begin{document}
\imprimirtitulodois
\begin{abstract}
Este ensaio examina as relações intersetoriais que compõem as diferenças entre as emissões atribuídas pela produção e pela demanda no Brasil. Combina a Matriz de Insumo-Produto de 2015, com 67 atividades, e intensidades de CO\textsubscript{2} em uma rede dirigida e ponderada. As forças e o número efetivo de contrapartes caracterizam volume e concentração; o filtro de disparidade seleciona relações conforme sua importância local. Apesar da conectividade quase uniforme, energia e água/resíduos distribuem volumes externos próximos com concentrações distintas, enquanto construção reúne emissões de diversas origens. As contrapartes esclarecem esses perfis e a importância assimétrica de uma mesma relação para origem e destino. O backbone retém 13,40\% das ligações e 70,69\% do peso intersetorial, com coberturas locais desiguais. A análise complementa a comparação das contas e informa relações a investigar para a coordenação da mitigação, com conclusões restritas ao cenário exploratório pelas limitações dos coeficientes.
\end{abstract}
\noindent\textbf{Palavras-chave:} atribuição de emissões; insumo-produto; redes dirigidas e ponderadas; mitigação.

\clearpage
\section{Introdução}

A atividade econômica está associada a externalidades ambientais, entre elas as emissões de dióxido de carbono (CO\textsubscript{2}), que contribuem para as mudanças climáticas. A redução das emissões integra os esforços globais de mitigação das mudanças climáticas \parencite{ipcc2023spm}. A alocação dos esforços de mitigação depende de quais agentes são responsabilizados por quais emissões. 

No âmbito do Protocolo de Kyoto, a contabilização das emissões segue a ótica da produção, onde as emissões são contabilizadas na atividade em que ocorrem \parencite{unfccc2008kyotoManual}. Existem outras óticas de atribuição setorial das emissões de CO\textsubscript{2}. Pela ótica da demanda, as emissões são atribuidas às atividades responsáveis pela demanda final, considerando o encadeamento das emissões diretas e indiretas \parencite{marques2012,millerBlair2009}. As emissões da geração de eletricidade, por exemplo, podem ser registradas no setor elétrico ou atribuídas aos produtos finais cuja produção requer essa eletricidade. A mudança de critério preserva o total contabilizado, mas altera sua distribuição entre as atividades.

A escolha de quem responsabilizar pelas emissões envolve uma dimensão normativa. A ótica da produção identifica a atividade em que a emissão ocorre, mas não determina, por si só, como distribuir os esforços de mitigação. Questões de justiça exigem reflexão e deliberação pública sobre os valores e os fins que devem orientar a vida em sociedade \parencite{sandel2015justica}. Para subsidiar essa deliberação, é necessário compreender como as diferentes óticas de atribuição distribuem as emissões entre as atividades produtivas.

A literatura sobre emissões brasileiras oferece elementos para comparar atribuições entre setores, regiões e cadeias produtivas \parencite{imori2015,sanguinetAzzoni2024,montoyaEtAl2026}. Os totais e os saldos setoriais, entretanto, não mostram, por si sós, quais produtos finais estão associados às emissões de cada atividade nem quais atividades emissoras compõem a atribuição de cada produto final. Também não mostram se os volumes atribuídos se concentram em poucas relações ou se distribuem entre muitas. Examinar essas relações intersetoriais permite complementar a comparação das contas pela produção e pela demanda.

Este trabalho pergunta: \textbf{Como as emissões de CO\textsubscript{2} se distribuem nas relações entre atividades emissoras e setores de produtos finais no Brasil em 2015?} O objetivo é caracterizar essas relações e discutir suas implicações para a coordenação dos esforços de mitigação.

A análise utiliza a extensão ambiental do modelo insumo-produto \parencite{millerBlair2009}, combinando a Matriz de Insumo-Produto (MIP) brasileira de 2015, com 67 atividades \parencite{ibge2015mip}, e coeficientes setoriais de emissão de \textcite{sanguinetAzzoni2024}. As atribuições são representadas em uma rede dirigida e ponderada, adaptando a abordagem de \textcite{huang2019} para relacionar as atividades emissoras aos setores de produtos finais.

O recorte abrange emissões da produção nacional associadas à demanda doméstica e às exportações, excluindo emissões estrangeiras. Os coeficientes ambientais compõem um cenário exploratório, com validação ambiental e monetária pendente, conforme detalhado na seção~\ref{sec:dados}.

No cenário exploratório, energia e água/resíduos apresentam volumes semelhantes de emissões atribuídas aos produtos finais de outras atividades, mas as atribuições de água/resíduos são mais concentradas. Os resultados também mostram que uma mesma relação pode ter importância diferente para a atividade emissora e para o setor do produto final.

\widowpenalty=10000
\clubpenalty=10000
\displaywidowpenalty=10000
% Sequência inspirada em Huang et al. (2019).
\clearpage
\input{dados.tex}
\clearpage
\section{Metodologia}
\input{modelo_atribuicao.tex}
\subsection{Construção da rede dirigida e ponderada}
\label{sec:construcao-rede}

A rede tem 67 nós, um por atividade da MIP, identificados pelos códigos do IBGE. Cada nó representa uma origem emissora nas ligações de saída e o setor de um produto final nas ligações de entrada, conforme a representação de \textcite{huang2019}. As categorias de uso final não constituem nós separados.

Para estudar relações entre atividades distintas, define-se
\begin{equation}
W_{ij}=\begin{cases}
F_{ij}, & i\ne j,\\
0, & i=j.
\end{cases}
\label{eq:rede-sem-diagonal}
\end{equation}
Cada $W_{ij}>0$ define uma aresta $i\rightarrow j$, cujo peso, em Gg de CO\textsubscript{2}, indica a emissão em $i$ atribuída ao produto final de $j$. As duas direções de um par são distintas. Como os pesos incorporam requisitos diretos e indiretos, concatenar arestas não representa uma sequência de etapas produtivas.

A diagonal $F_{ii}$ reúne emissões de $i$ atribuídas à demanda final de seus próprios produtos, incluindo requisitos indiretos. Ela permanece em $F$ e nas contas $c$ e $u$, mas é retirada de $W$, cujo peso total corresponde apenas às atribuições entre atividades distintas.

A construção preserva todos os nós e as arestas positivas, sem corte por magnitude. São conferidos o alinhamento dos códigos e a existência de pesos finitos e não negativos. A seleção de arestas é tratada na seção~\ref{sec:backbone}.

\subsection{Volume e concentração das relações}
\label{sec:volume-concentracao}

O volume e a concentração são calculados em $W$, antes da filtragem, por meio das forças e da disparidade dos pesos \parencite{huang2019,serrano2009}.

\paragraph{Volume.}
As forças de saída e entrada são as somas dos pesos em cada direção:
\begin{equation}
\begin{aligned}
s_i^{\mathrm{out}}&=\sum_{j\ne i}W_{ij}=c_i-F_{ii},\\
s_i^{\mathrm{in}}&=\sum_{j\ne i}W_{ji}=u_i-F_{ii}.
\end{aligned}
\label{eq:forcas-externas}
\end{equation}
A força de saída mede as emissões próprias de $i$ atribuídas aos produtos finais de outras atividades; a de entrada mede as emissões de outras atividades atribuídas aos produtos finais de $i$. A diferença entre força de entrada e força de saída é $s_i^{\mathrm{in}}-s_i^{\mathrm{out}}=u_i-c_i$. A leitura separada dos volumes revela atribuições que se compensam no saldo; as contas completas e a diagonal situam sua participação no total.

\paragraph{Concentração.}
Para perfis com força positiva, define-se a participação de cada contraparte:
\begin{equation}
p_{ij}^{\mathrm{out}}=\frac{W_{ij}}{s_i^{\mathrm{out}}},\qquad
p_{ji}^{\mathrm{in}}=\frac{W_{ji}}{s_i^{\mathrm{in}}},\quad j\ne i.
\label{eq:participacoes-externas}
\end{equation}
O perfil de saída distribui o volume intersetorial da origem entre destinos; o de entrada descreve as origens que compõem a atribuição intersetorial ao destino. Cada perfil soma um. A disparidade e seu inverso são
\begin{equation}
\begin{aligned}
Y_i^{\mathrm{out}}&=\sum_{j\ne i}(p_{ij}^{\mathrm{out}})^2,&
Y_i^{\mathrm{in}}&=\sum_{j\ne i}(p_{ji}^{\mathrm{in}})^2,\\
D_{2,i}^{\mathrm{out}}&=1/Y_i^{\mathrm{out}},&
D_{2,i}^{\mathrm{in}}&=1/Y_i^{\mathrm{in}}.
\end{aligned}
\label{eq:disparidade}
\end{equation}
$Y$ aumenta com a concentração; $D_2$ expressa a mesma informação como número efetivo de contrapartes. Para $k$ pesos iguais, $D_2=k$; quando uma relação domina, aproxima-se de um. As medidas são adimensionais e ficam indefinidas quando a força é zero, sem excluir o nó.

Os pesos e as participações identificam as contrapartes, cuja identidade não é informada por $D_2$. O grau conta as arestas positivas de entrada ou saída; a densidade é o número de arestas dividido por $67\times66$. Esses indicadores descrevem a presença de relações, enquanto a disparidade considera a distribuição dos pesos.

\subsection{Extração e avaliação do backbone}
\label{sec:backbone}

O filtro de disparidade de \textcite{serrano2009}, aplicado por \textcite{huang2019}, seleciona ligações destacadas na composição local dos pesos, considerando atividades de diferentes escalas. O backbone resultante não altera as contas ambientais completas.

\paragraph{Critério de seleção.}
O modelo nulo distribui a força entre $k$ ligações dividindo o intervalo unitário por $k-1$ pontos uniformemente escolhidos. Para $k>1$, a probabilidade de uma parcela atingir ou superar $p$ é $(1-p)^{k-1}$. Para a aresta $i\rightarrow j$, as caudas nas perspectivas de saída e entrada são
\begin{equation}
\begin{aligned}
a_{ij}^{\mathrm{out}}&=\left(1-\frac{W_{ij}}{s_i^{\mathrm{out}}}\right)^{k_i^{\mathrm{out}}-1},\\
a_{ij}^{\mathrm{in}}&=\left(1-\frac{W_{ij}}{s_j^{\mathrm{in}}}\right)^{k_j^{\mathrm{in}}-1}.
\end{aligned}
\label{eq:caudas-filtro}
\end{equation}
Os graus $k$ e as forças são os de $W$. A aresta é selecionada se $a_{ij}^{\mathrm{out}}<\alpha$ \textbf{ou} $a_{ij}^{\mathrm{in}}<\alpha$, registrando-se a extremidade responsável pela seleção. A saída destaca um destino nas atribuições da origem; a entrada destaca uma origem na composição do destino.\footnote{As caudas são calculadas apenas para graus pertinentes maiores que um. Se uma extremidade tem grau um, vale o teste da outra; se ambas têm grau um, a aresta é preservada por convenção de manutenção da relação exclusiva, inclusive em díades isoladas. Essa extensão dos casos degenerados do suplemento de Serrano et al. é registrada separadamente da seleção estatística.}

Adota-se $\alpha=0{,}05$, com $0{,}01$ e $0{,}10$ para sensibilidade. O filtro é aplicado uma vez a $W$ para cada valor, sem atualização dos perfis. A matriz $B_\alpha$ preserva os pesos originais nas arestas selecionadas e zero nas demais, sem renormalização ou corte absoluto adicional.

\paragraph{Avaliação da seleção.}
Comparam-se a fração de arestas retidas e a cobertura dos pesos, globalmente e por atividade:
\begin{equation}
\begin{aligned}
r_\alpha&=\frac{\sum_{ij}(B_\alpha)_{ij}}{\sum_{ij}W_{ij}},\\
r_{\alpha,i}^{\mathrm{out}}&=\frac{\sum_j(B_\alpha)_{ij}}{s_i^{\mathrm{out}}},\qquad
r_{\alpha,i}^{\mathrm{in}}=\frac{\sum_j(B_\alpha)_{ji}}{s_i^{\mathrm{in}}}.
\end{aligned}
\label{eq:cobertura-backbone}
\end{equation}
As razões são definidas apenas para denominadores positivos. A cobertura refere-se à parcela intersetorial de $W$ e permite distinguir retenção agregada de perdas por atividade e direção. A conectividade é descrita pelos graus, isolados e componentes fracamente conectados, obtidos ao desconsiderar a direção das arestas.

A sensibilidade é examinada pelas mudanças nas arestas, coberturas e posições setoriais entre os três valores de $\alpha$. O aninhamento dos conjuntos decorre da regra e não constitui evidência independente de robustez. As caudas do modelo nulo não medem incerteza ambiental, e a seleção não estabelece prioridade de mitigação.

Para interpretar os perfis, as contrapartes são ordenadas por peso em $W$ completa. Destacam-se as cinco maiores e a participação conjunta das demais, sem modificar a rede. A leitura dos casos combina essas participações, as contas completas e os motivos de retenção; a sensibilidade verifica também a permanência das principais relações nos três parâmetros.

\clearpage
\section{Resultados}
\label{sec:resultados}

Os resultados correspondem ao cenário ambiental exploratório descrito na seção~\ref{sec:dados}. A leitura parte da rede completa para distinguir volume e concentração das atribuições e identificar as contrapartes que compõem os perfis setoriais. Em seguida, examina quais relações o filtro destaca e quanto dessa composição permanece representado no backbone.

\subsection{Contas ambientais, volume e concentração}

As duas contas totalizam 683,70 mil Gg de CO\textsubscript{2}, dos quais 392,92 mil Gg estão na diagonal de $F$ e 290,78 mil Gg nas relações entre atividades distintas em $W$. A igualdade entre os totais pela produção e pela demanda decorre da identidade contábil. A parcela intersetorial permite examinar como as mesmas emissões conectam origens e produtos finais, enquanto a diagonal situa o volume atribuído aos produtos finais da própria atividade emissora.

A rede completa tem 67 nós e 4\,290 arestas, com densidade de 97,01\%. Serviços domésticos (9700) é o único isolado: sua conta permanece na diagonal de $F$. As outras 66 atividades estão ligadas entre si nas duas direções, com graus iguais a 65. A presença de relações é compatível com a incorporação dos requisitos indiretos em $F$ e pouco distingue os perfis setoriais. O conteúdo desses perfis precisa ser examinado pela distribuição dos pesos.

Entre os 66 perfis definidos, o número efetivo varia de 1,58 a 25,04 destinos nas saídas e de 2,13 a 22,10 origens nas entradas; as medianas são 11,50 e 13,93, respectivamente. A igualdade dos graus, portanto, convive com distribuições ponderadas distintas. A Figura~\ref{fig:indicadores-resultados} permite comparar volume, concentração e seleção nas atividades de maior força intersetorial.

Três atividades permitem desenvolver esses contrastes. Energia (3500) e água, esgoto e resíduos (3680) têm as duas maiores forças de saída, com volumes próximos de 31,32 e 31,23 mil Gg, mas apresentam 20,79 e 9,88 destinos efetivos. Construção (4180), a terceira maior força de entrada, recebe 23,48 mil Gg de outras origens e tem 20,76 origens efetivas. Os casos ilustram diferenças de concentração e de direção entre atividades de grande volume intersetorial; não constituem uma amostra representativa dos perfis.

As contas completas situam esses papéis (Tabela~\ref{tab:contas-energia-construcao}). Em energia e água/resíduos, a atribuição pela produção supera a atribuição pela demanda; em construção ocorre o inverso. No entanto, construção mantém 35,91 mil Gg na diagonal, equivalentes a 92,5\% de sua conta pela produção. Sua pequena saída intersetorial corresponde a uma parcela reduzida das emissões próprias, enquanto as entradas de outras origens ampliam sua conta pela demanda.

\clearpage
\newgeometry{margin=1.5cm}
\begin{landscape}
\begin{figure}[p]
\centering
\includegraphics[width=\linewidth,height=0.80\textheight,keepaspectratio]{figuras/indicadores.pdf}
\caption{Volume, concentração, conectividade selecionada e cobertura em atividades de maior volume externo.}
\label{fig:indicadores-resultados}
\begin{minipage}{\linewidth}\footnotesize
Fonte: elaboração própria. Linhas: união das seis maiores forças de saída e das seis maiores forças de entrada de $W$; não representam todas as atividades. Azul: saída; laranja: entrada. Volume e número efetivo são calculados em $W$ completa. Grau e cobertura referem-se a $\alpha=0{,}05$; os intervalos do grau abrangem $\alpha=0{,}01$--$0{,}10$ e não são intervalos de confiança. A diagonal é excluída.
\end{minipage}
\end{figure}
\end{landscape}
\restoregeometry

\begin{table}[!ht]
\centering
\caption{Contas e volumes intersetoriais das atividades examinadas.}
\label{tab:contas-energia-construcao}
\small
\begin{tabular}{@{}lrrr@{}}
\toprule
Medida & Energia & Água/resíduos & Construção \\
\midrule
Produção, $c$ & 56,62 & 51,80 & 38,84 \\
Demanda, $u$ & 28,23 & 21,74 & 59,40 \\
Atribuição à própria atividade, $F_{ii}$ & 25,30 & 20,57 & 35,91 \\
Saída externa, $s^{\mathrm{out}}$ & 31,32 & 31,23 & 2,93 \\
Entrada externa, $s^{\mathrm{in}}$ & 2,92 & 1,16 & 23,48 \\
Saldo, $u-c$ & $-28,39$ & $-30,07$ & 20,56 \\
\bottomrule
\end{tabular}
\par\smallskip\begin{minipage}{\textwidth}\footnotesize
Fonte: elaboração própria. Valores em mil Gg de CO\textsubscript{2}, antes do filtro. Códigos IBGE: energia (3500), água, esgoto e gestão de resíduos (3680) e construção (4180). A diagonal integra $c$ e $u$, mas não as forças externas. Diferenças residuais entre somas de valores exibidos decorrem do arredondamento.
\end{minipage}
\end{table}





\subsection{Composição dos perfis de atribuição}
\label{sec:composicao-perfis}

A composição dos destinos ajuda a interpretar a diferença de concentração (Tabela~\ref{tab:contrapartes-completas}). O maior destino das saídas de energia é comércio, com 13,49\% do volume externo, seguido de administração pública, construção, abate e produtos de carne e agricultura. Os cinco maiores destinos reúnem 40,45\% do volume, e os demais 60 respondem por 59,55\%. A distribuição envolve produtos finais de atividades industriais, agropecuárias e de serviços.

Em água/resíduos, administração pública concentra 27,22\% das saídas externas, seguida de comércio, saúde privada, saúde pública e construção. Os cinco maiores destinos reúnem 52,35\% do volume. Assim, a maior concentração decorre, em particular, do peso da atribuição à administração pública, cuja participação é mais de três vezes a observada nas saídas de energia. Os destinos identificam atividades dos produtos finais, conforme a fronteira do modelo; não representam categorias institucionais de compradores nem apenas aquisições diretas de energia ou de serviços de água e resíduos.

Na perspectiva de entrada, construção reúne emissões de diferentes origens. Produtos de minerais não metálicos representam 12,60\% do volume externo recebido, seguidos de energia, água/resíduos, comércio e produtos de metal. As cinco maiores origens somam 38,71\%, deixando 61,29\% distribuídos entre outras 60 atividades. Sua conta pela demanda combina, portanto, as emissões atribuídas aos produtos finais da própria construção, na diagonal, com emissões intersetoriais distribuídas entre diversas origens. A posição da construção nas duas contas é esclarecida pela composição desses volumes, e não apenas pelo saldo positivo.

\clearpage
\begin{table}[!ht]
\centering
\caption{Principais contrapartes nos perfis completos, antes do filtro.}
\label{tab:contrapartes-completas}
\small
\begin{tabular}{@{}>{\raggedright\arraybackslash}p{0.65\textwidth}rr@{}}
\toprule
Contraparte (código IBGE) & Peso & Participação (\%) \\
\midrule
\multicolumn{3}{@{}l}{\textit{Saídas de energia (3500): produtos finais}} \\
Comércio (4580) & 4,224 & 13,49 \\
Administração pública, defesa e seguridade social (8400) & 2,556 & 8,16 \\
Construção (4180) & 2,263 & 7,23 \\
Abate e produtos de carne, laticínio e pesca (1091) & 1,930 & 6,16 \\
Agricultura (0191) & 1,693 & 5,41 \\
Demais 60 destinos & 18,650 & 59,55 \\
\midrule
\multicolumn{3}{@{}l}{\textit{Saídas de água/resíduos (3680): produtos finais}} \\
Administração pública, defesa e seguridade social (8400) & 8,501 & 27,22 \\
Comércio (4580) & 2,964 & 9,49 \\
Saúde privada (8692) & 1,672 & 5,35 \\
Saúde pública (8691) & 1,621 & 5,19 \\
Construção (4180) & 1,592 & 5,10 \\
Demais 60 destinos & 14,880 & 47,65 \\
\midrule
\multicolumn{3}{@{}l}{\textit{Entradas de construção (4180): origens emissoras}} \\
Produtos de minerais não metálicos (2300) & 2,959 & 12,60 \\
Energia elétrica, gás natural e outras utilidades (3500) & 2,263 & 9,64 \\
Água, esgoto e gestão de resíduos (3680) & 1,592 & 6,78 \\
Comércio (4580) & 1,240 & 5,28 \\
Produtos de metal, exceto máquinas e equipamentos (2500) & 1,037 & 4,42 \\
Demais 60 origens & 14,393 & 61,29 \\
\bottomrule
\end{tabular}
\par\smallskip\begin{minipage}{\textwidth}\footnotesize
Fonte: elaboração própria. Peso em mil Gg de CO\textsubscript{2}. As participações usam a força externa da atividade e direção indicadas, sem diagonal. As cinco maiores contrapartes são uma seleção de exposição, sem alterar $W$ ou aplicar corte adicional. Nomes abreviados; valores calculados sem arredondamento no notebook. A relação energia $\rightarrow$ construção aparece nas duas perspectivas e não deve ser somada duas vezes em totais da rede.
\end{minipage}
\end{table}

As direções inversas também distinguem os dois casos. Energia recebe 2,92 mil Gg de outras origens, com 21,59 origens efetivas; construção distribui 2,93 mil Gg a outros produtos finais, com 8,10 destinos efetivos. Neste último perfil, administração pública reúne 31,21\% do volume externo. Mesmo com forças próximas, a distribuição das saídas de construção é mais concentrada do que a das entradas de energia. Cada nó precisa, assim, ser interpretado nos dois papéis que exerce na rede.

A composição dos perfis esclarece a passagem entre as contas. Em energia e água/resíduos, as saídas para outros produtos finais superam as entradas de outras origens, reduzindo a atribuição pela demanda em relação à produção. Em construção, as entradas superam as saídas e elevam essa atribuição. A diagonal integra as duas contas e cancela na diferença; as contrapartes identificadas mostram quais relações intersetoriais compõem os saldos.

\clearpage
\subsection{Seleção das relações e limites de cobertura}

O backbone de referência preserva 575 arestas, ou 13,40\% das ligações de $W$, e 70,69\% do peso intersetorial (Tabela~\ref{tab:sensibilidade-backbone}). Na Figura~\ref{fig:matriz-backbone}, a linha de energia e a coluna de construção mostram a extensão das relações retidas nesses dois perfis. O painel de pesos diferencia os volumes dessas ligações: a seleção por importância local não implica pesos absolutos semelhantes.

\begin{table}[htbp]
\centering
\caption{Sensibilidade da seleção e da cobertura ao parâmetro do filtro.}
\label{tab:sensibilidade-backbone}
\small
\begin{tabular}{@{}crrrr@{}}
\toprule
$\alpha$ & Arestas & \% das arestas de $W$ & Peso retido & Cobertura (\%) \\
\midrule
0,01 & 338 & 7,88 & 166,88 & 57,39 \\
0,05 & 575 & 13,40 & 205,55 & 70,69 \\
0,10 & 755 & 17,60 & 224,92 & 77,35 \\
\bottomrule
\end{tabular}
\par\smallskip\begin{minipage}{\textwidth}\footnotesize
Fonte: elaboração própria. Peso em mil Gg de CO\textsubscript{2}. Cobertura relativa ao peso de $W$, sem diagonal. Todos os casos mantêm 67 nós, um isolado e um componente fracamente conectado de 66 atividades. O filtro é aplicado à mesma rede completa, com regra OU e sem atualização dos perfis.
\end{minipage}
\end{table}

A cobertura agregada oculta perdas locais. Em $\alpha=0{,}05$, as menores coberturas são 32,96\% nas saídas e 33,79\% nas entradas. Portanto, preservar parcela expressiva do peso total não assegura representação igualmente completa de cada atividade e direção.

A mudança de $\alpha$ altera as posições individuais por grau, embora o componente de 66 atividades permaneça nos três casos. A conectividade agregada não demonstra estabilidade das posições setoriais. Os valores detalhados das mudanças, a decomposição da seleção por extremidade e a sensibilidade dos perfis focais constam do \ref{anexo:diagnosticos}; $0{,}05$ permanece como referência, sem ser considerado um parâmetro ótimo.





\clearpage
\newgeometry{margin=1.5cm}
\begin{landscape}
\begin{figure}[p]
\centering
\includegraphics[width=\linewidth,height=0.82\textheight,keepaspectratio]{figuras/matriz_presenca.pdf}
\caption{Matriz de adjacência do backbone com $\alpha=0{,}05$: (a) presença das arestas.}
\label{fig:matriz-backbone}
\begin{minipage}{\linewidth}\footnotesize
Fonte: elaboração própria. Linha: origem emissora; coluna: atividade do produto final. Ordem setorial do IBGE nos dois eixos. Azul identifica as 575 ligações retidas; cinza indica ausência no backbone, inclusive a diagonal. Ausência no backbone não implica valor nulo em $F$ ou $W$. Todas as 67 atividades permanecem nos eixos.
\end{minipage}
\end{figure}
\clearpage
\begin{figure}[p]
\ContinuedFloat
\centering
\includegraphics[width=\linewidth,height=0.82\textheight,keepaspectratio]{figuras/matriz_peso.pdf}
\caption{Matriz de adjacência do backbone com $\alpha=0{,}05$: (b) pesos das mesmas arestas.}
\begin{minipage}{\linewidth}\footnotesize
Fonte: elaboração própria. Mesma ordem e mesma seleção do painel (a). A cor representa pesos originais em Gg de CO\textsubscript{2}, em escala logarítmica, sem normalização por linha ou coluna. Cinza identifica células sem aresta retida. A cor mede volume absoluto; a concentração é descrita pelos números efetivos em $W$ completa.
\end{minipage}
\end{figure}
\end{landscape}
\restoregeometry

\subsection{Importância local nas duas perspectivas da relação}

Os perfis completos mostram como os pesos se distribuem; o filtro acrescenta a avaliação da importância local de cada relação para suas duas extremidades. As cores da Figura~\ref{fig:rede-focos} evidenciam o predomínio de saídas no recorte de energia e de entradas no de construção, mantendo as demais relações do mesmo backbone ao fundo. O motivo da retenção e a cobertura, apresentados na Tabela~\ref{tab:motivos-casos} e nos painéis do \ref{anexo:paineis}, permitem interpretar essa diferença entre as duas vistas.

\begin{table}[htbp]
\centering
\caption{Motivos de retenção e cobertura nos recortes de energia e construção.}
\label{tab:motivos-casos}
\small
\begin{tabular}{@{}lrrrrr@{}}
\toprule
Recorte & Arestas & Só saída & Só entrada & Ambos & Cobertura (\%) \\
\midrule
Energia: entradas & 7 & 1 & 5 & 1 & 42,1 \\
Energia: saídas & 62 & 0 & 56 & 6 & 99,2 \\
Construção: entradas & 45 & 41 & 0 & 4 & 91,4 \\
Construção: saídas & 5 & 4 & 1 & 0 & 57,0 \\
\bottomrule
\end{tabular}
\par\smallskip\begin{minipage}{\textwidth}\footnotesize
Fonte: elaboração própria. $\alpha=0{,}05$. Saída e entrada nas colunas de motivo identificam o teste na origem e no destino da aresta, independentemente da atividade em destaque. Cobertura: peso retido dividido pela força da mesma atividade e direção em $W$. Os recortes compartilham arestas e não devem ser somados para obter totais da rede.
\end{minipage}
\end{table}

Em energia, 56 das 62 saídas retidas passam somente pelo teste de entrada no destino. Energia tem participação destacada na composição externa de muitas demandas finais, embora cada destino possa representar parcela pequena das emissões distribuídas pela origem. Por exemplo, energia $\rightarrow$ alimentação (5600) representa 3,00\% das saídas externas de energia e 4,86\% das entradas externas de alimentação. A relação é selecionada somente na perspectiva do destino, evidenciando a diferença entre a composição recebida por uma atividade e a distribuição das emissões pela outra.

Em construção, 41 das 45 entradas passam somente pelo teste de saída na origem. Construção representa parcela destacada das atribuições externas dessas origens, mesmo quando a contribuição individual é pequena no total recebido pelo destino. Produtos de metal (2500) $\rightarrow$ construção, por exemplo, corresponde a 32,57\% das saídas externas da origem e a 4,42\% das entradas externas do destino. O filtro retém essa ligação somente pelo teste de saída na origem. Ela identifica uma demanda final que concentra parte expressiva das atribuições externas de produtos de metal, mesmo ocupando uma posição menor na composição recebida por construção.

As entradas de construção selecionadas pelos dois testes vêm de produtos de minerais não metálicos, energia, água/resíduos e comércio. A relação energia $\rightarrow$ construção, comum aos dois recortes, pesa aproximadamente 2\,263 Gg: 7,23\% das saídas externas de energia e 9,64\% das entradas externas de construção. Sua seleção nas duas extremidades conecta os papéis predominantes dos dois casos.

As saídas de energia e as entradas de construção têm as maiores coberturas dos respectivos recortes (Tabela~\ref{tab:motivos-casos}). Nas direções inversas, as listas curtas não autorizam desprezar as relações omitidas em conjunto. Tampouco a diferença entre 62 saídas de energia e 45 entradas de construção indica concentração distinta: na rede completa, os respectivos números efetivos são próximos, ambos de 20,8 contrapartes.

A sensibilidade qualifica a leitura desses papéis. Entre $\alpha=0{,}01$ e $0{,}10$, energia mantém de 54 a 64 saídas, com cobertura de 92,12\% a 99,96\%, enquanto construção mantém de 32 a 55 entradas, com cobertura de 74,31\% a 97,76\% (Tabela~\ref{tab:sensibilidade-casos}). Em todos os parâmetros, predomina o teste somente de entrada no destino nas saídas de energia e somente de saída na origem nas entradas de construção. Os cinco maiores destinos de energia permanecem selecionados. Esses resultados sustentam a leitura dos papéis ilustrados dentro da faixa examinada, apesar das mudanças no número de ligações e na cobertura.

A permanência de cada relação, contudo, exige verificação própria. Entre as cinco maiores origens de construção, água/resíduos deixa de ser selecionada em $\alpha=0{,}01$, embora represente 6,78\% do perfil externo recebido. Energia $\rightarrow$ construção permanece selecionada pelos dois testes nos três parâmetros. A composição da rede completa e os saldos não mudam com $\alpha$; a instabilidade de uma ligação no backbone delimita sua seleção sob o modelo nulo, sem alterar o volume de sua atribuição.



\clearpage
\begin{figure}[p]
\centering
\includegraphics[width=\textwidth,height=0.76\textheight,keepaspectratio]{figuras/rede_foco_3500.pdf}
\caption{Duas perspectivas do mesmo backbone: (a) destaque para energia (3500).}
\label{fig:rede-focos}
\begin{minipage}{\textwidth}\footnotesize
Fonte: elaboração própria. Cada vista mantém 67 nós e 575 arestas em $\alpha=0{,}05$. Azul: saídas do destaque; laranja: entradas; cinza: demais arestas. Energia tem 62 saídas e sete entradas, com cobertura de 99,2\% e 42,1\% do volume externo de cada direção. Área dos nós uniforme; espessura cresce com a raiz quadrada do peso. Posição definida pelo sinal de $u-c$ e pela ordem IBGE: 33 nós com $c>u$ à esquerda, 33 com $u>c$ à direita e o isolado abaixo. A disposição não resulta de centralidade ou comunidade.
\end{minipage}
\end{figure}
\clearpage
\begin{figure}[p]
\ContinuedFloat
\centering
\includegraphics[width=\textwidth,height=0.76\textheight,keepaspectratio]{figuras/rede_foco_4180.pdf}
\caption{Duas perspectivas do mesmo backbone: (b) destaque para construção (4180).}
\begin{minipage}{\textwidth}\footnotesize
Fonte: elaboração própria. Mesmos nós, arestas, posições e escala de espessura do painel (a). Construção tem cinco saídas e 45 entradas, com cobertura de 57,0\% e 91,4\%. Cores representam direção em relação ao destaque, não motivo de retenção. As duas vistas preservam 70,69\% do peso intersetorial de $W$. Cruzamentos e proximidade não representam trajetórias ou distâncias econômicas; relações individuais podem ser consultadas na matriz e nos painéis do anexo.
\end{minipage}
\end{figure}
\clearpage

\clearpage
\section{Discussão e implicações para a coordenação da mitigação}

A diferença entre as contas pela produção e pela demanda informa o saldo de uma atividade, mas não identifica as contrapartes nem a concentração dos volumes que o compõem. Energia e água/resíduos atribuem a outros produtos finais volumes superiores aos provenientes de outras origens; em construção, predomina o sentido inverso. A leitura conjunta das contas e dos perfis esclarece essas diferenças: a diagonal mantém a parcela atribuída aos produtos finais da própria atividade, e as relações intersetoriais identificam as emissões atribuídas a setores distintos de sua origem. Essa leitura complementa as comparações de atribuição entre setores, regiões e cadeias produtivas presentes na literatura brasileira \parencite{imori2015,sanguinetAzzoni2024,montoyaEtAl2026}.

Volumes semelhantes também podem envolver conjuntos de contrapartes com participações diferentes. A maior concentração das saídas de água/resíduos, em comparação com energia, destaca o peso da administração pública no primeiro perfil. Já a concentração das saídas de construção não caracteriza suas entradas, mais distribuídas entre origens. Esses contrastes mostram por que a posição de uma atividade nas contas não basta para descrever sua participação nas relações intersetoriais. A representação dirigida de \textcite{huang2019} permite distinguir os dois papéis de cada atividade, enquanto a análise dos pesos considera diferenças que o grau quase uniforme da rede completa não informa.

A importância local distingue, ainda, relações destacadas para a origem ou para o destino. Construção pode representar uma parcela expressiva das atribuições externas de produtos de metal sem que essa origem domine a composição recebida pela construção. Assim, examinar apenas os maiores componentes da conta de um destino pode deixar em segundo plano relações importantes na perspectiva das origens. O filtro de \textcite{serrano2009} explicita essa assimetria sob seu modelo nulo. A seleção precisa ser lida com as participações e as coberturas, pois as ligações omitidas podem representar, em conjunto, parcela expressiva de um perfil.

Em uma aplicação a Santa Catarina em 2018, esse procedimento permitiria complementar a comparação das participações setoriais com duas questões: quais produtos finais estão associados às emissões dos setores cuja participação diminui na atribuição pela demanda e quais origens compõem a atribuição dos setores cuja participação aumenta. A concentração distinguiria perfis com volumes mais equilibrados entre contrapartes daqueles dominados por poucas relações. A aplicação exigirá dados harmonizados e o recálculo das contas e dos perfis na fronteira estadual; as contrapartes e os contrastes do cenário nacional não podem ser transportados para o estado.

Para a coordenação da mitigação, origens compartilhadas por várias demandas e demandas associadas a várias origens indicam conjuntos de atividades a investigar conjuntamente. No primeiro caso, cabe examinar como mudanças em uma origem emissora se relacionariam com os produtos finais aos quais suas emissões são atribuídas. No segundo, cabe investigar o uso de energia e materiais associado a um produto final considerando suas diferentes origens emissoras. Intervenções nos processos produtivos e no uso de energia e materiais ao longo das cadeias constituem frentes complementares \parencite{rissman2020}. Os resultados ajudam a formular essas questões, mas não estimam o efeito de mudanças tecnológicas ou de demanda, a capacidade de intervenção dos agentes, os custos ou a eficácia de instrumentos. A distribuição normativa dos esforços de mitigação requer critérios adicionais à atribuição contábil.

As interpretações ambientais permanecem condicionadas à base utilizada. A interpolação, a cobertura não reconciliada e a diferença de preços limitam a aplicação dos contrastes setoriais à economia brasileira observada. Essas restrições alcançam os volumes e as comparações entre atividades, podendo também afetar os perfis ponderados e a seleção de relações. A consistência contábil e a sensibilidade a $\alpha$ verificam aspectos distintos da validade das emissões físicas. Antes de sustentar conclusões ambientais definitivas, é necessário reconciliar a fronteira, a alocação setorial e os denominadores monetários dos coeficientes ou substituí-los por uma extensão compatível com 2015.

\clearpage
\section{Conclusões}

No cenário analisado, as emissões de energia e água/resíduos atribuídas a outros produtos finais superam as emissões provenientes de outras origens e atribuídas aos produtos finais de energia e água/resíduos. Em construção ocorre o inverso, ampliando sua conta pela demanda em relação à produção. A identificação das contrapartes esclarece esses saldos e distingue perfis de volumes próximos com concentrações diferentes, apesar da conectividade semelhante. Construção também mostra a necessidade de separar as emissões atribuídas aos produtos finais da própria atividade das relações intersetoriais. A leitura conjunta das contas e dos perfis permite compreender a composição das diferenças entre produção e demanda.

O backbone distingue relações destacadas para a origem, para o destino ou para ambos. Os casos mostram que essa importância pode ser assimétrica e que a estabilidade dos papéis examinados convive com mudanças na seleção e na cobertura. A análise de redes oferece, assim, um procedimento para aprofundar a comparação das participações setoriais e formular questões de coordenação da mitigação. Seu aproveitamento em Santa Catarina requer uma aplicação própria, e conclusões empíricas definitivas sobre o Brasil dependem da validação ou substituição da extensão ambiental.

\clearpage
{\renewcommand*{\bibfont}{\small}
\printbibliography[title={Referências}]
}
\input{anexos.tex}
\end{document}
